\chapter{Conclusion, Limitations \& Future Work}
\label{chap:conclusion}

\section{Summary of Achievements}
Throughout this course project for Data Mining (CO3029), team \textbf{hí ộ} successfully designed, implemented, and empirically validated an advanced personalized literature recommendation system:

\begin{enumerate}
  \item \textbf{Curricular \& Methodological Mastery:}
    \begin{itemize}
      \item Substantively synthesized core Data Mining methods: semantic \textit{K-Means Clustering}, frequent itemset mining via \textit{Apriori}, and multi-model gradient boosting ensembles (\textit{LightGBM, XGBoost, CatBoost}).
      \item Formulated a robust \textbf{Two-Stage Recommendation Pipeline} uniting high-precision binary pruning with a novel \textit{Multiclass Expected Gain} ranking objective.
      \item Analytically established the theoretical score ceiling ($\approx 0.776$) via Monte Carlo simulations, demystifying the trade-off mechanics governing the official evaluation metric.
    \end{itemize}

  \item \textbf{Empirical Leaderboard Impact:}
    \begin{itemize}
      \item Progressively boosted model performance from baseline ($0.61546$) to the Sentence-BERT ensemble ($0.65006$), securing **Rank 5 on the Kaggle Leaderboard** within only two submission iterations.
      \item Maintained strict 5-Fold Stratified Cross-Validation protocols with zero data leakage, ensuring competitive endurance on the final Private Test set.
    \end{itemize}

  \item \textbf{Engineering Rigor \& Open Science:}
    \begin{itemize}
      \item Automated the complete end-to-end pipeline from text normalization to submission generation within a structured GitHub repository: \url{https://github.com/sangngng/DataMining261}.
    \end{itemize}
\end{enumerate}

\section{Current Limitations}
Despite strong competitive results, several intrinsic constraints remain:
\begin{enumerate}
  \item \textbf{Absence of Abstracts and Full Text:} The dataset provides solely paper titles. Academic titles are necessarily succinct and occasionally fail to convey methodological scope, occasionally precipitating false positive predictions for broad survey papers.
  \item \textbf{Co-Authorship Network Sparsity:} Over 4,200 unique authors appear in the collection, the majority having only a single publication. Modeling preferences for newly observed first-time authors remains constrained by data sparsity.
  \item \textbf{Moderate Sample Volume:} A training base of 2,494 papers is relatively modest for complex representation learning, making pre-trained transformer embeddings essential.
\end{enumerate}

\section{Future Research Directions}
To push recommendation quality toward the theoretical ceiling, we propose:
\begin{enumerate}
  \item \textbf{Graph Neural Networks (GNNs) on Heterogeneous Citation Networks:} Constructing a heterogeneous graph linking authors, venues, and publications. Graph Convolutional Networks (GCN) or Graph Attention Networks (GAT) could capture transitive community influence and collaborative filtering signals far beyond flat feature tables.
  \item \textbf{Large Language Model (LLM) Reasoning \& Re-Ranking:} Utilizing advanced reasoning models (e.g., LLaMA-3, Mistral) to read candidate titles in context, provide explainable recommendations, and dynamically re-rank the top 30 candidate publications.
  \item \textbf{Automated Metadata Enrichment via Academic APIs:} Leveraging CrossRef and Semantic Scholar APIs to automatically ingest paper abstracts, citation counts, and author affiliations based on registered DOIs, providing enriched signals for downstream rankers.
\end{enumerate}
